\documentclass[12pt]{article}
\usepackage{amsmath}
\usepackage{amssymb}
\usepackage{geometry}
\geometry{a4paper, margin=1in}

\title{Lecture Scribe \\ Gaussian, Uniform, Exponential, Gamma Random Variables}
\author{}
\date{}

\begin{document}

\maketitle

\section*{Run 1 : Scribe with Reasoning}

\subsection*{1. Lecture Topic Identification}

\textbf{Reasoning:}  
From first slide, topic is Gaussian, Uniform, Exponential, Gamma Random Variables.  
Title slide defines scope of lecture.

\textbf{Final:}  
Lecture covers continuous random variables with focus on Gaussian distribution.

\hrule

\subsection*{2. Outline Understanding}

\textbf{Reasoning:}  
Slides show Gaussian → Continuous RV types → Examples → Problem solving.  
Outline defines order of concepts.

\textbf{Final:}

Lecture flow:

\begin{itemize}
\item Gaussian RV
\item Uniform / Exponential / Laplace / Gamma
\item Examples
\item CDF analysis
\item Applications
\end{itemize}

\hrule

\subsection*{3. Moments of Random Variable}

\textbf{Reasoning:}  
Board notes show moments used to describe distribution.

\textbf{Final:}

Mean
\[
\mu = E[X]
\]

Variance
\[
\sigma^2 = E[(X-\mu)^2]
\]

Skewness : measure of symmetry

Kurtosis : measure of peak

\hrule

\subsection*{4. Expectation of Function}

\textbf{Reasoning:}  
Used later in Gaussian analysis.

\textbf{Final:}

For Continuous RV

\[
E[g(X)] = \int g(x) f_X(x) dx
\]

For Discrete RV

\[
E[g(X)] = \sum g(x_k) p_k
\]

Linearity property

\[
E[aX + b] = aE[X] + b
\]

\hrule

\subsection*{5. Gaussian Random Variable Definition}

\textbf{Reasoning:}  
Slide defines Gaussian PDF.

\textbf{Final:}

\[
f_X(x) =
\frac{1}{\sqrt{2\pi\sigma^2}}
\exp \left(
-\frac{(x-m)^2}{2\sigma^2}
\right)
\]

Gaussian depends on mean $m$ and variance $\sigma^2$.

\hrule

\subsection*{6. Standard Forms}

\textbf{Reasoning:}  
Used for probability calculations.

\textbf{Final:}

Error function : erf(x)

Complementary error : erfc(x)

\[
\Phi(x) = \text{CDF}
\]

\[
Q(x) = 1 - \Phi(x)
\]

\hrule

\subsection*{7. Relation between $\Phi$ and $Q$}

\textbf{Reasoning:}  
Derived using standardization.

\textbf{Final:}

\[
F_X(x) =
\Phi \left(
\frac{x-m}{\sigma}
\right)
\]

\[
Pr(X > x) =
Q \left(
\frac{x-m}{\sigma}
\right)
\]

CDF uses $\Phi$  
Right tail uses $Q$

\hrule

\subsection*{8. Example Gaussian Problem}

Given

\[
m = -3
\quad
\sigma = 2
\]

\textbf{Final:}

\[
Pr(X \le 0) =
\Phi \left(
\frac{3}{2}
\right)
\]

\[
Pr(X > 4) =
Q \left(
\frac{7}{2}
\right)
\]

\[
Pr(|X+3| < 2)
\]

\[
Pr(|X-2| > 1)
\]

\hrule

\subsection*{9. Applications of Gaussian}

\textbf{Reasoning:}  
Slides show engineering examples.

\textbf{Final:}

\begin{itemize}
\item Thermal noise
\item Measurement error
\item Network delay
\end{itemize}

Gaussian models uncertainty.

\hrule

\subsection*{10. CDF Analysis Problem}

Given

\[
F_X(x) =
\begin{cases}
0, & x < 0 \\
x^2, & 0 \le x \le 1 \\
1, & x > 1
\end{cases}
\]

\textbf{Reasoning:}  
Differentiate to get PDF.

\textbf{Final:}

\[
f_X(x) = 2x
\]

Mean

\[
E[X] = \frac{2}{3}
\]

Variance

\[
Var(X) = \frac{1}{18}
\]

Skewness < 0

Kurtosis < Gaussian

\hrule

\subsection*{11. Gaussian Modeling}

\textbf{Reasoning:}  
Standard normal used in modeling.

\textbf{Final:}

\[
X = \sigma Z + \mu
\]

Gaussian describes uncertainty.

\hrule

\subsection*{12. Density Estimation}

\textbf{Reasoning:}  
Samples used to estimate parameters.

\textbf{Final:}

Observed samples → estimate $\mu$, $\sigma$

Statistics gives estimated PDF.

\hrule

\subsection*{13. Applications in Computing}

\textbf{Reasoning:}  
Slides show modern applications.

\textbf{Final:}

\begin{itemize}
\item Image processing
\item Cloud latency
\item IoT sensors
\end{itemize}

Probability used in modern engineering.

\end{document}